\iffalse
\section{DISCUSION}
La limitacion del experimento realizado es que solo podemos obtener el valor de
$\epsilon \cdot \sigma$. Con lo cual, si queremos usar el resultado para estimar
alguno de estos parametros, es necesario obtener el valor del otro a partir de
otra experiencia.

\subsection*{Calculo de sigma a partir de una emisividad tabulada}
ver comentario en el codigo
% amigo no le veo nada de sentido a hacer este.
% Ese alguien calculo ese epsilon en funcion del sigma tabulado y sus mediciones.
% Si nuestro resultado tiene error, calcular algo con ese valor va a tener doblemente error.
% Calculemos el epsilon y lo comparemos con el epsilon de ese alguien y veamos por que da tanto mas chico.
% No se cual decis vos.

\subsection*{Calculo de la emisividad a partir de sigma tabulado}
Tomando como referencia el valor tabulado
$\sigma = 5.670374 \cdot 10^{-8}\,\mathrm{W\,m^{-2}\,K^{-4}}$, los valores de
emisividad $\epsilon$ inferidos se resumen en la tabla \ref{tab:epsilon}.

\begin{table}
    \begin{center}
        \begin{tabular}{| c | c |}
            relacion & $\epsilon$ \\ \hline
            C1 - C2 & $0.65 \pm 0.09$ \\
            C2 - C3 & $0.735 \pm 0.007$ \\ \hline
        \end{tabular}
        \caption{Valores de $\epsilon$ obtenidos a partir de los ajustes realizados.}
        \label{tab:epsilon}
    \end{center}
\end{table}

Estos valores son demasiado bajos y no son comparables con el valor nominal
0.97 asignado a la pintura negra mate en informes previos de Gatto y
Z\'u\~niga\cite{gatto}. Al combinar en cuadratura la incerteza
instrumental con la del ajuste, las estimaciones de $\sigma$ para las relaciones
$C2$-$C3$ quedan limitadas en buena medida por la propagacion instrumental,
mientras que en $C1$-$C2$ sigue dominando la dispersion del ajuste. La
diferencia en los valores que involucran al cilindro 1, en particular la
mostrada en el grafico \ref{fig:fit12}, puede deberse a una falla en el
termometro utilizado para el cilindro 1 o en algun punto del circuito. Sin
embargo, la relacion entre los cilindros 2 y 3 tampoco da un valor de
emisividad cercano, asi que tambien es posible que se haya da\~nado la pintura de
los cilindros, que los cilindros no sean equivalentes o que algun valor usado
en el procedimiento no haya sido bien medido.
\fi

\section{DISCUSION}
